\ProvidesPackage{cncs-pseudocode}

\RequirePackage{algpseudocode}
\RequirePackage{tcolorbox}

% ---------------------------------------------------------------------------
% Thai / English keyword sets for algpseudocode.
%
% algpseudocode does NOT have separate "\algorithmicendif", "\algorithmicendfor"
% macros — \EndIf, \EndFor, etc. are auto-built by \algdef as
% "\algorithmicend + \algorithmicif" (see algpseudocode.sty). So the only
% macros that need redefining are the base keyword atoms below; "end if",
% "end for", ... follow automatically once \algorithmicend and the matching
% keyword are both redefined.
%
% \cncs@setthaikeywords / \cncs@setenglishkeywords are only ever invoked
% inside a TeX group (see \pseudocode / \pseudocodethai below), so the
% redefinition never leaks into a later block in the same document.
% ---------------------------------------------------------------------------
\newcommand{\cncs@setthaikeywords}{%
  \algrenewcommand\algorithmicend{\textbf{จบ}}%
  \algrenewcommand\algorithmicdo{\textbf{ทำ}}%
  \algrenewcommand\algorithmicwhile{\textbf{ขณะที่}}%
  \algrenewcommand\algorithmicfor{\textbf{สำหรับ}}%
  \algrenewcommand\algorithmicforall{\textbf{สำหรับทุก}}%
  \algrenewcommand\algorithmicloop{\textbf{วนซ้ำ}}%
  \algrenewcommand\algorithmicrepeat{\textbf{ทำซ้ำ}}%
  \algrenewcommand\algorithmicuntil{\textbf{จนกระทั่ง}}%
  \algrenewcommand\algorithmicprocedure{\textbf{กระบวนการ}}%
  \algrenewcommand\algorithmicfunction{\textbf{ฟังก์ชัน}}%
  \algrenewcommand\algorithmicif{\textbf{ถ้า}}%
  \algrenewcommand\algorithmicthen{\textbf{แล้ว}}%
  \algrenewcommand\algorithmicelse{\textbf{มิฉะนั้น}}%
  \algrenewcommand\algorithmicrequire{\textbf{นำเข้า:}}%
  \algrenewcommand\algorithmicensure{\textbf{ผลลัพธ์:}}%
  \algrenewcommand\algorithmicreturn{\textbf{คืนค่า}}%
}

% Restates algpseudocode's own English defaults explicitly, so the two
% keyword sets sit side by side and stay easy to compare/edit together.
\newcommand{\cncs@setenglishkeywords}{%
  \algrenewcommand\algorithmicend{\textbf{end}}%
  \algrenewcommand\algorithmicdo{\textbf{do}}%
  \algrenewcommand\algorithmicwhile{\textbf{while}}%
  \algrenewcommand\algorithmicfor{\textbf{for}}%
  \algrenewcommand\algorithmicforall{\textbf{for all}}%
  \algrenewcommand\algorithmicloop{\textbf{loop}}%
  \algrenewcommand\algorithmicrepeat{\textbf{repeat}}%
  \algrenewcommand\algorithmicuntil{\textbf{until}}%
  \algrenewcommand\algorithmicprocedure{\textbf{procedure}}%
  \algrenewcommand\algorithmicfunction{\textbf{function}}%
  \algrenewcommand\algorithmicif{\textbf{if}}%
  \algrenewcommand\algorithmicthen{\textbf{then}}%
  \algrenewcommand\algorithmicelse{\textbf{else}}%
  \algrenewcommand\algorithmicrequire{\textbf{Require:}}%
  \algrenewcommand\algorithmicensure{\textbf{Ensure:}}%
  \algrenewcommand\algorithmicreturn{\textbf{return}}%
}

% ---------------------------------------------------------------------------
% Shared themed box (same visual language as \cncscode: primary-colour frame,
% cncsbg background). \begin{algorithmic}[1] turns on line numbers, matching
% cncscode's default numbered look.
% ---------------------------------------------------------------------------
\newtcolorbox{cncs@pseudobox}{%
  colback=cncsbg!50,
  colframe=cncsprimary,
  boxrule=1pt,
  arc=1mm,
  left=8pt, right=8pt, top=4pt, bottom=4pt,
  breakable
}

% \begin{pseudocode} ... \end{pseudocode}
% Keywords in English (if/then/else/for/while/...); body text, variable
% names, and \Comment{...} can still be typed in Thai.
\newenvironment{pseudocode}{%
  \begingroup
  \cncs@setenglishkeywords
  \cncs@pseudobox
  \begin{algorithmic}[1]
}{%
  \end{algorithmic}
  \endcncs@pseudobox
  \endgroup
}

% \begin{pseudocodethai} ... \end{pseudocodethai}
% Same box and numbering, but every control keyword is in Thai.
\newenvironment{pseudocodethai}{%
  \begingroup
  \cncs@setthaikeywords
  \cncs@pseudobox
  \begin{algorithmic}[1]
}{%
  \end{algorithmic}
  \endcncs@pseudobox
  \endgroup
}

% ---------------------------------------------------------------------------
% \begin{pseudosteps} ... \end{pseudosteps} — Thai-native numbered-outline
% style, matching how Thai textbooks actually write pseudocode: nesting is
% shown by decimal step numbers (4, 4.1, 4.2, ...), NOT by indentation +
% "end if"/"end for" keywords the way algpseudocode (\pseudocode above) does.
% There is no block-closing keyword at all in this style — a block simply
% ends when the numbering returns to the parent level.
%
%   \begin{pseudosteps}
%     \item ให้ $n$ \pcgets จำนวนข้อมูลในรายการ $L$
%     \item ให้ $left$ \pcgets $1$
%     \item ทำซ้ำ ขณะที่ $left \le right$
%       \begin{pseudosubsteps}
%         \item ให้ $mid$ \pcgets $(left+right)/2$ ปัดเศษทิ้ง
%         \item ถ้า $x = target$ แล้ว ให้คืนค่าดัชนี $i$ เท่ากับ $mid$ และจบการทำงาน
%       \end{pseudosubsteps}
%     \item คืนคำตอบว่าไม่พบข้อมูล $target$ ในรายการ $L$
%   \end{pseudosteps}
%
% \pcgets prints the assignment arrow (←) via math mode \leftarrow rather
% than the literal Unicode glyph, so it never depends on the Thai font
% having that character (the \texttt-in-Thai bug earlier this session was
% exactly this kind of font/glyph mismatch — this sidesteps it entirely).
%
% \PCBegin / \PCEnd print a plain "เริ่มต้น" / "จบ" heading — OPTIONAL,
% not part of the environment itself: some worked examples use them to
% bookend the whole algorithm, some don't (both patterns appear in real
% textbook material), so it's the author's call per problem, not forced.
% Place them OUTSIDE \begin{pseudosteps}...\end{pseudosteps}, not inside —
% a LaTeX list environment must start with \item, so \PCBegin{} (which
% prints a \par) breaks the list machinery if it's the first thing inside:
%
%   \PCBegin
%   \begin{pseudosteps}
%     \item ...
%   \end{pseudosteps}
%   \PCEnd
% ---------------------------------------------------------------------------
\RequirePackage{enumitem}

\newcommand{\pcgets}{\ensuremath{\leftarrow}}
\newcommand{\PCBegin}{\par\noindent\textbf{เริ่มต้น}\par}
\newcommand{\PCEnd}{\par\noindent\textbf{จบ}\par}

\newenvironment{pseudosteps}{%
  \cncs@pseudobox
  \begin{enumerate}[label=\arabic*., leftmargin=1.8em, itemsep=0.3em, topsep=0.2em]
}{%
  \end{enumerate}
  \endcncs@pseudobox
}

% Must be used INSIDE a \pseudosteps \item (i.e. genuinely nested), so it
% picks up the parent step's number automatically via enumitem's own
% depth-tracking (enumi at depth 1, enumii at depth 2) — no manual number
% bookkeeping, so renumbering a step never desyncs the sub-step labels.
\newenvironment{pseudosubsteps}{%
  \begin{enumerate}[label=\arabic{enumi}.\arabic*, leftmargin=2.8em, itemsep=0.2em, topsep=0.2em]
}{%
  \end{enumerate}
}

% Third nesting level (3.1.1-style) — must be used INSIDE a \pseudosubsteps
% \item, same depth-tracking mechanism (now enumi.enumii.self at depth 3).
% Not needed often — keep nesting shallow where you can, it's easier for
% students to follow — but the mechanism costs nothing extra to support
% when a lesson genuinely needs it.
\newenvironment{pseudosubsubsteps}{%
  \begin{enumerate}[label=\arabic{enumi}.\arabic{enumii}.\arabic*, leftmargin=3.8em, itemsep=0.15em, topsep=0.15em]
}{%
  \end{enumerate}
}
